\chapter{Operation}

\section{Command line}
\index{command line}\index{mock mode}\index{operation}%
\begin{table}[h]
  \centering
  \begin{tabularx}{\textwidth}{lD}
    \toprule
    \textbf{Option} & \textbf{Effect} \\
    \midrule
    \cfg{-m}, \cfg{--mock} & Use simulated robot, conveyor and camera. No
      hardware required. \\
    \cfg{-c}, \cfg{--config <path>} & Configuration file
      (default \cfg{Settings.toml}). \\
    \cfg{-h}, \cfg{--help} & Show help. \\
    \bottomrule
  \end{tabularx}
  \caption{Command-line options}
\end{table}

\index{logging}\index{RUST LOG@\texttt{RUST\_LOG}}%
Logging uses \cfg{env\_logger}; set \cfg{RUST\_LOG} to control verbosity:
\begin{lstlisting}[language=bash]
RUST_LOG=info cargo run -- --mock      # normal operation
RUST_LOG=debug cargo run               # includes G-code level tracing
\end{lstlisting}

\section{Startup sequence}
\index{startup sequence}\index{handshake}\index{homing}%
On start the application: loads and validates the configuration; connects
to the robot (\cfg{IsDelta} handshake, Appendix~\ref{app:gcode}) and homes
it if \cfg{home\_on\_connect} is set; connects to the conveyor; acquires
the emergency-stop handles; connects to the camera; starts the
orchestrator; starts the vision loop (which takes ownership of the camera
and turns detected objects into pick requests); and finally opens the
operator UI. Any failure in this chain aborts startup with a message
identifying the failing device.

\section{Automatic sorting}
\index{sorting}\index{vision}\index{pick}%
While the belt runs, the vision loop detects objects, follows each one
across several frames and hands the orchestrator exactly one pick request
per object, with the position converted to robot coordinates using the
calibration of section~\ref{sec:calibration}. Two safeguards apply on the
receiving side: pick requests are discarded while the orchestrator is
paused (including after an emergency stop or a command failure), and any
pick request older than 3~seconds is discarded as stale, since the belt
will have carried the object away. Discarded picks are logged as
warnings. If frame capture fails, the vision loop logs a warning and
retries every 500~ms without affecting robot control.

\section{Operator interface}
\index{user interface}\index{controls}\index{dashboard}%
\begin{table}[h]
  \centering
  \begin{tabularx}{\textwidth}{lD}
    \toprule
    \textbf{Control} & \textbf{Behaviour} \\
    \midrule
    Start Sorting / Pause System & Toggles the conveyor. Starts the belt at
      \cfg{default\_speed}; pressing again stops it. The label reflects the
      current state. \\
    Stop Conveyor & Stops the belt (enabled while running). \\
    Home Robot & Queues a homing command through the orchestrator (enabled
      while the system is not running). \\
    EMERGENCY STOP & Immediately halts robot and conveyor and clears all
      queued work. See section~\ref{sec:estop}. \\
    \bottomrule
  \end{tabularx}
  \caption{UI controls}
\end{table}

The dashboard shows the robot connection status, the conveyor state and
the total number of sorted bricks. The live camera feed and detection
overlays are not wired up yet (see the TODO list).

\section{Emergency stop and recovery}
\label{sec:estop}
\index{emergency stop}\index{recovery}\index{M112@\texttt{M112}}%
Pressing \textbf{EMERGENCY STOP}:
\begin{enumerate}
  \item writes \cfg{M112} to the robot and \cfg{M5} to the conveyor on
    dedicated serial handles, bypassing every queue and lock;
  \item clears the orchestrator's instruction queue and pauses it;
  \item marks both devices as \cfg{E-STOP} in the UI.
\end{enumerate}

To recover after an E-stop:
\begin{enumerate}
  \item Remove the physical cause (jam, obstruction, misplaced fixture).
  \item Restart the application. The robot controller may also need a
    power cycle after \cfg{M112}, depending on firmware.
  \item Home the robot before resuming production.
\end{enumerate}

\begin{notebox}
  If a robot command fails during operation (timeout, serial error, robot
  fault), the orchestrator clears its queue and pauses itself
  automatically rather than continuing with a desynchronised state. The
  incident is logged with the failing command.
\end{notebox}

\section{Routine operation checklist}
\index{checklist}%
\begin{enumerate}
  \item Check that the belt is empty and the workspace is clear.
  \item Start the application; confirm \emph{Status: Connected}.
  \item If not homed automatically, press \emph{Home Robot}.
  \item Press \emph{Start Sorting}.
  \item Supervise the first picks; keep a hand near the E-stop.
\end{enumerate}
